\section{Verification}

Six kinds of check were made, each testing something further from the code than the one before
(Table~\ref{tab:verification}). Identity and reproduction checks test the code alone: a patch at its default
leaves the arithmetic of the build without it unchanged along the trajectory tested, and the recorded code and
results can be rebuilt. Arithmetic checks test an implementation against a formula or a direct solution. A
similarity check tests whether gravity enters the code consistently, by showing that a run at scaled gravity,
lengths and times reproduces the original in dimensionless form. Calibrations and benchmarks compare with
something outside the code: the reference model on the columns a fit was made on, and a scheme's published
behaviour on Earth. Each leaves the lunar climate untested in its own way, since the reference is itself a
clear-sky model resting on its spectroscopy, a fit holds near its fitted state, similarity leaves Earth's
empirical laws untested at lunar gravity, and a benchmark runs on Earth.

\textbf{Identity and reproduction.} Five ExoPlaSim runs, each a model year from one restart, repeat byte for
byte the restart file of the same build without the patch, with the patch at its neutral setting: the
Rayleigh and continuum calibration at the released values, the cloud-water and convective-day patches,
infrared from space at zero, and the ice optics off. One of them compares against the released model itself,
and the other four against builds that already had earlier patches, so the chain back to the release closes
through those patches' own identities; a year killed a minute in and resumed repeats the uninterrupted year's
restart and every variable of its output.\claim{V01} In CM1 the tilted-ring patches and the local-time
advection repeat an existing run's first model day bit for bit, single-thread runs repeat, and the electrified
supercell, restarted on four threads and on two, repeats its fields and its 41 flashes; SWAN with the
wind-input patch writes stock output at standard gravity.\claim{V02} Applied to the pinned upstream archives,
the patch series rebuild byte for byte the deployed ExoPlaSim installation and the Fortran of the eight lunar
CM1 builds used for the results, and prefixes of the ExoPlaSim series reproduce the source hashes of 32 of
the 34 GCM runs that recorded one.\claim{V03} The Sun clock, the water scaling, and the curved-air, gas-band
and longwave patches (07--09) were written to reproduce the released code at their off settings, and have yet
to be run that way.

\textbf{Similarity.} The strongest evidence about gravity comes from the similarity checks. Stock SWAN, run at
a sixth of Earth's gravity with lengths, depths and durations six times longer and frequencies six times
lower, misses its Earth twin by 5.44\% in significant wave height and 3.91\% in mean period; patched, it misses
by 0.303\% and 0.326\%, against a 1\% criterion.\claim{V04} SWASH reproduces the linear dispersion relation
within 1.76\% in amplitude and 0.289\% in phase at lunar, four times lunar and Earth gravity, preserves a linear
waveform to 0.0009\% of its amplitude when gravity quadruples and the clock halves, and, under the same scaling
with breaking and a moving shoreline, changes its 98th-percentile run-up by 1.20\%.\claim{V05} Each lunar wave
run is therefore the gravity-scaled image of an Earth run with the same Earth-fitted coefficients, and whether
those coefficients hold on lunar seas is untested. CM1's gravity-scaled microphysics has no such test yet:
the pair of runs designed to measure how far lunar convection departs from Earth's stretched sixfold failed to
convect.\claim{V06}

\textbf{Calibration and the reference.} The calibrations hold on the columns they were fitted to and a little
beyond. The clear-sky split of sunlight holds within 0.6 points rms on eleven columns outside the fit's set,
and a 303~K column carrying 338~kg\,m$^{-2}$ of vapour, warmer and wetter than any fitted, falls 0.7 points
short, where the scheme before its water recalibration fell 9.7 short.\claim{V07} RRTMG, checked and left
uncalibrated, lies within 2.5 points of the line-by-line shares of reflected, air-absorbed and ground-absorbed
sunlight at the site where the rings cross the equator at 0\textdegree~E, and on 990 cloud-free columns of the
45\textdegree\,N ring cools the air more than the correlated-$k$ model does, by 4.5~W\,m$^{-2}$ in the lowest
1.8~km, 8.5 at 1.8--6.5~km and 10.2 at 6.5--13.7~km.\claim{V08} Every calibration is clear-sky, and PlaSim's
clouds remain uncalibrated. The accuracy of all of them is bounded by the reference's own. Its correlated-$k$
emulation reproduces the line-by-line outgoing longwave within 0.4--1.4~W\,m$^{-2}$, and the line-by-line
model has two checks against other codes: PyRADS, for longwave with water vapour alone at constant gravity,
and a Monte Carlo reference for Rayleigh scattering up to an optical depth of 1.5, where the lunar column
reaches about 4 at 360~nm. No independent code has yet checked its carbon dioxide, oxygen, collision-induced
absorption or sunlight on the lunar column.\claim{V16}

\textbf{Benchmarks and arithmetic.} The electrified graft was benchmarked on Earth. Against CM1's own NSSL
scheme in the Weisman--Klemp supercell \autocite{weisman1982dependence}, its largest updrafts stay within 6\%,
its reflectivity within 3.4~dBZ at every half hour and its condensate within 9\%, while hail differs by up to
a factor of two, the two copies being different versions of the scheme.\claim{V09} Against published runs of
the same charging and lightning schemes, the first flash comes at 28.5 minutes against a published 27--40, and
the flash rates and the median charge per flash fall within the published ranges; with WRF-ELEC's ground rule
on, 1.1\% of flashes struck the ground, against 0.33\% in a published run. The charging runs several times
stronger, traced to cloud ice reaching $2\times10^8$ crystals per kilogram where droplets freeze homogeneously
in 50~m\,s$^{-1}$ updrafts.\claim{V10} The lunar conductivity column reproduces Earth's measured fair-weather
positive conductivity within 2--8\% from 5 to 30~km at a 3~GV cutoff. Arithmetic checks rebuild the uneven
sunlight of the Sun-clock defect from PlaSim's solar-position arithmetic with correlation 1.0000, place a
tilted ring's Sun within $8\times10^{-7}$ of the spherical formula and its Coriolis parameter within
$3\times10^{-12}$~s$^{-1}$ of $2\Omega\sin\phi$, match the field solver to a direct solve to one part in
$10^5$, and hold the curved-air ledger's channels within 0.9~W\,m$^{-2}$ of the fine spectral grid in every
layer with the Sun overhead.\claim{V11}

\textbf{What the checks caught.} The adaptation introduced defects of its own, found by these checks or in
the course of the work. The GCM runner first extended the solar spectrum past its 2.73~\textmu m end at its
last value, an invented infrared tail that put 86\% of the sunlight in the near infrared and warmed the first
run past 320~K.\claim{H28} The first electrified build discharged once per 6~s step, after all the step's
charging, so the field stood a median 3.2 times breakdown at a step's first discharge, which then removed a
median 150~C. Its fall-speed scaling multiplied memory the scheme had left unset, and an OpenMP directive made
private a variable the threads shared, both caught when results changed with the thread count; and the first
lunar gravity patch scaled the NSSL fall speeds but left the ventilation tied to them at Earth's values, so
lunar rain evaporated 1.3--1.5 times too fast for its size until a code review found the omission.\claim{H29}
The spherical sky solver's production grid of 114 spectral channels averaged optical depth over windows and
opaque lines and absorbed 58\% too much sunlight above 500~nm with the Sun overhead, which had made PlaSim look
short of its calibration; the ledger's first channels, bisected on the whole column alone, put near-infrared
sunlight too high, 19.0 against 10.7~W\,m$^{-2}$ at 66--85~km, and the first two 20-year GCM tests of the
curved air ran on those terms.\claim{H30} The first terrain build placed its forcing by a corner column
standing 2.9~km up; CM1's starting soundings held the vapour of the GCM's top level, near 130~km, constant
above it, so the 45\textdegree\,N rings began 6--15 times saturated over ice at 140--150~km and every column
kept 1.1~g\,m$^{-2}$ of ice above 107~km through both lunar days; and the 45\textdegree\,N rings with and
without the curved-air terms, the twilight pair, ran on builds compiled at different optimisation levels, so
their 0.013--0.021~K difference lies below what the pair resolves.\claim{H31} The last two remain in the runs
made before they were found.

\textbf{Checks that fell short.} SWAN with depth breaking off failed its repeatability test by about 2.5\%, and
tracing the failure found the uninitialised breaker parameter. SWAN's early growth depends on the time step:
halving it changes the first hours' wave height by up to 33\%, while the 48-hour endpoints move under 0.3\%,
and small seas in the weather runs fail their 5\% time-step criterion where seas above 1~m pass. SWASH's control
with irregular input misses its 5\% bound, at 5.22\% in height and 6.28\% in period.\claim{V12} CM1 rebuilt at a
higher optimisation departs from its first build by 0.28--0.32~\textdegree C rms in the air at 2~m within a
model day, inside the spread of last-bit nudges, and on several threads it departs from itself. The ledger's
scattering orders reached their cap in one channel of 282, whose omitted tail is about $3\times10^{-4}$ of that
channel's diffuse light at the ground.\claim{V13}

\textbf{Checks still to make.} The full design column has yet to be run through an independent radiation code
with every absorber, which would test the reference on which all the calibrations rest. CM1 has yet to be
forced from the longwave-corrected GCM, so whether the corrected physics closes the gap between the two
models is untested. The longwave correction, fitted at 295~K, ran in a test that settled near 307~K with twice
the water vapour, and its accuracy there is unchecked. The twilight pair lacks a control on the same build and
an estimate of its realisation noise. CM1's thread race remains unlocated, RRTMG's extrapolation past its
tables and CM1's gravity-scaled microphysics remain untested on their own, and an independent GCM has yet to
run the scenario.\claim{V14}
